\scanpage{25}
% archive sheet number: 24
$\mathbb C$, on a un objet canonique \deleted{de} $\unsure{\xi|\zeta}$ de $T_{\mathbb C}(\mathbb Z)$ [$T$ le torseur fondamental associé à $\mathbb C$).
\[
 H^i(X,\mathbb Q(0))\rightsquigarrow H^i(X_{\mathrm{top}},\mathbb Z)
\]
L'hom $H^i(X^n,\mathbb Q(0))\longrightarrow H^{i}(Y^{\unsure{m|''}},\mathbb Q(0))$ défini par \deleted{$\unsure{A^n}$} $z\in A^n(X\times Y)$ est celui défini, \inserted{par voie \unsure{transc.}, à l'aide de} \deleted{par} la classe $\unsure{P|\rho}_{\mathrm{top}}(z)\left[=\frac{1}{(2i\pi)^n}P_{\mathrm{alg}}(z)\right]$
\[
 \boxed{\xi(z)(x)=\mathrm{pr}_{2*}^{\mathrm{top}}\bigl(\mathrm{pr}_1^*(x)P_{\mathrm{top}}(z)\bigr)}
\]
Considérons l'\deleted{\unsure{isom}} can (pour $X,i$ donnés)
\diagram{\begin{tikzcd}[ampersand replacement=\&,column sep=small,row sep=large]
 {\xi_{\mathbb C}(H^{i}(X,\mathbb Q(0))=\xi(H^i(X,\mathbb Q(0))\otimes_{\mathbb Z}\mathbb C}\arrow[r,"\sim"] \&
 {\underline{\underline H}^{\deleted{\unsure{*}}i}(X,\unsure{\Omega|\underline{\Omega}}_X^*)=\xi_{\mathbb C,\infty}(H^i(X,\mathbb Q(0))}\\
 {H^{i}(X_{\unsure{\mathrm{top}}},\mathbb C)}\arrow[u,phantom,"{\shortparallel}" description]\arrow[ur,dash,"{\sim\;\deleted{\unsure{\mathrm{can}|\mathrm{car}}}}","{\substack{\mathrm{can\ par}\\\text{Hodge-Gaga}}}"'] \&
\end{tikzcd}}

Cet isom \uline{\deleted{\illegible{1}} est \deleted{\illegible{1}} fonctoriel} en $H^{\unsure{p|i}}(X,\mathbb Q(0))$.

En effet, \unsure{si} $z\in A^n(X\times Y)$, on a \deleted{Pour le rendre canonique, il faut}
\[
\begin{aligned}
 x\ \rightsquigarrow\ \mathrm{pr}_{2*}^{\mathrm{top}}\bigl(\mathrm{pr}_1^*(x)P_{\mathrm{top}}(z)\bigr)
 &=\mathrm{pr}_{2*}^{\mathrm{top}}\left(\mathrm{pr}_1^{\unsure{a|*|\mathrm{an}}}(x)\left(\frac{1}{2i\pi}\right)^n P_{\mathrm{alg}}(z)\right)\\
 &=\left(\frac{1}{2i\pi}\right)^{-n}\mathrm{pr}_{2*}^{\mathrm{alg}}\left(\mathrm{pr}_1^*(x)\left(\frac{1}{2i\pi}\right)^n P_{\mathrm{alg}}(z)\right)
\end{aligned}
\]
\diagram{\begin{tikzcd}[ampersand replacement=\&,column sep=huge,row sep=huge]
 {H^{\unsure{p|i}}(X,\mathbb C)}\arrow[r,"{\xi_{\mathbb C}(z)\deleted{\unsure{\mathrm{alg}}}}","{\substack{\text{défini par}\\ P_{\mathrm{top}}(z)\\=(\frac{1}{2i\pi})^nP_{\mathrm{alg}}(z)}}"']\arrow[d,dash,"{\wr\;\mathrm{can}}"] \&
 {H^{\unsure{p|i'}}(Y,\mathbb C)}\arrow[d,dash,"{\wr\;\mathrm{can}}"]\\
 {\underline{\underline H}^p(X,\Omega_X^\cdot)}\arrow[r,"{\substack{\xi_{\mathbb C\infty}^{\unsure{\mathrm{de}}}(z)\\\text{défini par }P_{\mathrm{alg}}(z)\\=(\frac{1}{2i\pi})^n\xi_{\mathbb C}(z)}}"' ] \&
 {H^{\unsure{p|i'}}(Y,\underline\Omega_{\unsure{x|X|Y}}^\cdot)}
\end{tikzcd}}
% slanted
\margin{Cette compatibilité peut déjà \unsure{s'énoncer} pour analytique $\unsure{\longleftrightarrow|\longrightarrow}$ \unsure{topologique} (avec \deleted{\unsure{classes|chgt}} de corr. analytiques)}
\[
 x\ \rightsquigarrow\ \mathrm{pr}_{2*}^{\mathrm{alg}}\bigl(\mathrm{pr}_1^*(x)\underbrace{P_{\mathrm{alg}}(z)}\bigr)
\]
\hfill Ouf !

